\section{Implementación y verificación}
El contorno simplificado tiene 2061 vértices. Se resuelve $\Delta(h-\ln R)=0$ con borde $\ln\cos\varphi$, mediante diferencias finitas y Shortley--Weller junto a la frontera \cite{shortley1938}. Se usa $\Delta w=0.001$, con 58312 nodos. Las aristas proporcionan incrementos de conjugada y mapa, conciliados por mínimos cuadrados. Un ajuste polinomial previo dejó residuos cercanos a $9\times10^{-4}$, por lo que no se eligió para esta geometría \cite{brubeck2021}.

Con la solución conocida $h=\ln|w-w_*|$, singularidad exterior, los errores máximos son $1.46\times10^{-4}$, $3.74\times10^{-5}$, $9.74\times10^{-6}$ y $2.49\times10^{-6}$ con pasos $0.004$, $0.002$, $0.001$ y $0.0005$: convergencia aproximadamente cuadrática. Las mallas $0.002$, $0.001$ y $0.0005$ dan $\delta=0.00184343$, $0.00184347$ y $0.00184298$. En cabeceras, la prueba a 1 km conserva el ángulo recto a $0.003^\circ$. En los vértices de frontera, la oscilación interpolada de $\ln\sigma$ es $1.11\times10^{-4}$ y la desviación máxima respecto a $c$ es $7.52\times10^{-5}$. Estos residuos reflejan la reconstrucción discreta. La evaluación local puede tener saltos entre celdas de unos 6 mm: se documentan como error numérico, no como propiedad conforme exacta.

Mercator y Lambert esféricas coinciden con PROJ a menos de $10^{-8}$ m \cite{proj}. Los parámetros corresponden a EPSG:6372, \emph{Mexico ITRF2008 / LCC}, no a EPSG:6361, ausente como CRS en la base consultada \cite{epsg}. El código conserva funciones puras para fórmulas y clases para malla y modelo persistente. La carga común evita resolver de nuevo durante cada análisis.
